\documentclass[a4paper,11pt]{article}

\usepackage[french]{babel}
\usepackage[utf8]{inputenc}
\usepackage[T1]{fontenc}
\usepackage{amsmath,amssymb}
\usepackage{geometry}

\geometry{margin=2.5cm}

\begin{document}

\title{Produit matriciel}
\date{}
\maketitle

\section*{Base d'entraînement}

\bigskip

\subsection*{Exemple 1}
Petites dimensions pour découvrir le produit matriciel.

\vspace{2cm}

\[
\begin{array}{cc}
\, &

\left(
\begin{array}{cc}
-3 & 2 \\
1  & 1 
\end{array}
\right) \\

\left(
\begin{array}{cc}
1 & 4
\end{array}
\right) &

\left(
\begin{array}{cc}
A & F
\end{array}
\right)
\end{array}
\]

\subsection*{Exemple 2}
Une matrice qui transforme un mot en un autre, fait sens aux calculs

\vspace{2cm}

\[
\begin{array}{cc}
\, &

\left(
\begin{array}{ccc}
0 & -1 & 2 \\
1  & -2 & 4 \\
1 & 1 & 0 
\end{array}
\right) \\

\left(
\begin{array}{ccc}
5 & 1 & 21
\end{array}
\right) &

\left(
\begin{array}{ccc}
V & I & N
\end{array}
\right)
\end{array}
\]

\section*{Enigme}

\[
\begin{array}{cc}
\, &

\left(
\begin{array}{rrrr}
10 & -1 & -5 & 0 \\
0  & 0  & 0  & 0 \\
0  & 0  & 0  & 0 \\
0  & 1  & 0  & 1 \\
-3 & 0  & 4  & -1
\end{array}
\right) \\

\left(
\begin{array}{ccccc}
3 & 12 & 15 & 19 & 5
\end{array}
\right) &

\left(
\begin{array}{cccc}
\phantom{-0} & \phantom{-0} & \phantom{-0} & \phantom{-0}
\end{array}
\right)
\end{array}
\]

\subsection*{Réponse}
OPEN

\end{document}